% !TEX root = ../../main.tex
\subsection{安全模型创新}

FoodShield 采用的是“日常匿名、异常可溯源”的可控匿名安全模型。该模型不同于完全实名模式，也不同于完全匿名模式。完全实名模式虽然便于平台追责，但会增加用户真实身份、联系方式和订单偏好暴露的风险；完全匿名模式虽然保护隐私，但在发生恶意骚扰、虚假投诉或配送纠纷时又可能导致责任主体难以确认。FoodShield 在两者之间建立了一个受控平衡点。

在日常通信状态下，系统通过 PID 匿名身份、订单级通信房间和字段最小化展示机制，限制骑手端和管理员端能够看到的信息范围。骑手端不展示用户真实身份，不展示用户 PID，不展示用户端私有备注和标签；管理员端默认不展示通信明文，而是以订单摘要、消息哈希、Merkle Root 和审计日志作为主要查看对象。这种设计降低了普通业务流程中的身份暴露风险。

在异常处理状态下，系统允许用户端和骑手端基于订单号提交溯源申请，管理员经过身份认证后处理申请，并将相关操作写入审计日志。溯源入口从 PID 调整为订单号，使溯源流程更符合真实外卖平台中的投诉和纠纷处理逻辑，也降低了 PID 被直接滥用查询的风险。管理员的查询、验证和溯源行为均被纳入审计范围，从而避免审计权限本身成为新的隐私风险。

因此，FoodShield 的安全模型创新不只是“隐藏身份”，而是将匿名性、认证性、完整性、可追责性和可审计性统一到同一套订单生命周期中，实现了隐私保护与平台治理之间的动态平衡。
